\chapter{The Second Output}
% full version: chapters/16_chemoutput.tex

The machine has two outputs. The fast one is the muscles. The slow one is the body's chemistry: the brain controls the heart, the blood vessels, and the glands, and sends messages through the blood to the whole body at once. The blood is the most broadcast bus of all: one message reaches every cell in about a minute, and, as always, its meaning is given by the recipient's ``lock.''

\section*{Gas and Brake}

The heart beats by itself; it has its own metronome. The brain only adjusts the tempo with two wires of opposite sign. \nt{NORA} is the gas pedal: it speeds the heart up, but slowly. \nt{ACET} is the brake: it slows it down fast, because it does not need a long chain of chemical reactions. And here at last the \nt{ADRE} type finds a job: adrenaline released into the blood is the same gas pedal, only pressed for the whole body at once.

\pencilfig{16}{\pencilart{16}}{The heart is driven by two wires of opposite sign: gas and brake.}

\section*{A Cascade with Feedback}

Many hormones are controlled by a cascade: the brain orders the first gland, that gland orders the second, the second secretes a hormone, and the hormone tells the brain there is already enough of it. This is a three-stage thermostat. An engineer knows what is dangerous about such a circuit: if the loop is made too ``sensitive,'' it starts to swing, and the level no longer holds but oscillates. So the precision of such a regulator is limited: it cannot be made both stable and very precise at the same time.

\pencilfig{127}{%
% three identical first-order stages with proportional feedback: secretion = max(0, K (target - level)).
% K = 3 settles at 3/4 of the target; K = 12 (above the stability limit K = 8) keeps oscillating. x = 0.42 t, y = 0.55 level
\foreach \yo/\t in {1.75/{weak loop: calm, but off target}, 0/{strong loop: nearer the target, but it swings}}{
  \draw[pen] (0,\yo) -- (8.5,\yo);
  \draw[pcGray, densely dashed, line width=0.5pt] (0,\yo+0.55) -- (8.5,\yo+0.55);
  \node[lbl, anchor=south west] at (2.0,\yo+0.8) {\t};}
\node[lbl, anchor=west] at (8.55,2.3) {target}; \node[lbl, anchor=west] at (8.55,0.55) {target};
\begin{scope}[yshift=1.75cm]
\draw[penlite=pcBlue] plot[smooth] coordinates {(0.00,0.00) (0.17,0.01) (0.34,0.08) (0.50,0.19) (0.67,0.33) (0.84,0.46) (1.01,0.55) (1.18,0.59) (1.34,0.58) (1.51,0.55) (1.68,0.49) (1.85,0.43) (2.02,0.37) (2.18,0.34) (2.35,0.33) (2.52,0.34) (2.69,0.37) (2.86,0.40) (3.02,0.43) (3.19,0.44) (3.36,0.45) (3.53,0.45) (3.70,0.44) (3.86,0.42) (4.03,0.41) (4.20,0.40) (4.37,0.39) (4.54,0.39) (4.70,0.40) (4.87,0.40) (5.04,0.41) (5.38,0.42) (5.88,0.42) (6.38,0.41) (7.06,0.41) (8.40,0.41)};
\end{scope}
\draw[penlite=pcRed] plot[smooth] coordinates {(0.00,0.00) (0.17,0.05) (0.34,0.30) (0.50,0.68) (0.67,1.02) (0.84,1.23) (1.01,1.30) (1.18,1.26) (1.34,1.16) (1.51,1.02) (1.68,0.87) (1.85,0.72) (2.02,0.58) (2.18,0.47) (2.35,0.39) (2.52,0.38) (2.69,0.46) (2.86,0.57) (3.02,0.67) (3.19,0.71) (3.36,0.70) (3.53,0.65) (3.70,0.58) (3.86,0.50) (4.03,0.43) (4.20,0.41) (4.37,0.45) (4.54,0.53) (4.70,0.61) (4.87,0.65) (5.04,0.64) (5.21,0.60) (5.38,0.54) (5.54,0.47) (5.71,0.42) (5.88,0.42) (6.05,0.48) (6.22,0.56) (6.38,0.62) (6.55,0.64) (6.72,0.62) (6.89,0.56) (7.06,0.50) (7.22,0.44) (7.39,0.42) (7.56,0.45) (7.73,0.53) (7.90,0.60) (8.06,0.63) (8.23,0.62) (8.40,0.58)};
\node[lbl, anchor=north] at (4.2,-0.05) {time};
}{A model of a three-stage hormone cascade. Weak feedback settles the level but leaves it well below the target. Strong feedback brings it closer, but then the level no longer holds; it swings.}

\section*{Pulses, Not a Level}

Some hormones work only in pulses. If they are supplied constantly, the receiving gland ``tires'' and shuts down; if in short pulses once an hour, it works. The recipient reads not the level but the rhythm, like the synapse from the Morse code chapter that hears only changes.

\section*{The Day}

Finally, the body has a slow oscillator with a period of about a day. Its natural period differs slightly from twenty-four hours, and every morning light tunes it, the way a radio receiver tunes to a station. Do not confuse it with a computer's clock generator: it does not measure out the steps of a computation. It switches the operating modes of the whole machine: day and night, waking and sleep.

\fullversion{16}
